\chapter{Experimental Methodology}
\label{chap:methodology}

This chapter describes the experimental design, datasets, and evaluation
methodology used in this DDP.  The experiments are organised into two
groups: four preliminary validation experiments (EV1--EV4) conducted on
a simplified single-SKU topology, and six main optimisation experiments
(E0--E4) on the full 5-SKU 1N3 network.

%% -----------------------------------------------------------------------
\section{Preliminary Validation Experiments (EV1--EV4)}
\label{sec:meth:ev}

Before running the main optimisation suite, four controlled experiments
were designed to verify that the simulator correctly reproduces known
theoretical behaviour.  These experiments used a simplified \textbf{single-SKU,
2-warehouse $\times$ 3-retailer} topology and real M5 Walmart demand data.

\begin{table}[htbp]
  \centering
  \caption{Preliminary validation experiments: scope and purpose.}
  \label{tab:ev-overview}
  \begin{tabular}{@{}llp{7cm}@{}}
    \toprule
    ID & Title & Purpose \\
    \midrule
    EV1 & Constant-demand convergence &
          Verify that base-stock policy drives on-hand inventory to the
          target level under deterministic constant demand. \\[4pt]
    EV2 & Demand shock response &
          Verify that a one-time spike propagates upstream with pipeline
          delay and amplification. \\[4pt]
    EV3 & Policy comparison under transport capacity &
          Compare base-stock vs $(s,S)$ with and without vehicle capacity
          limits; confirm batching effect. \\[4pt]
    EV4 & Capacity scaling on M5 data &
          Quantify the non-linear cost--service trade-off as transport
          capacity is scaled from 0.25$\times$ to 4$\times$ baseline. \\
    \bottomrule
  \end{tabular}
\end{table}

\subsection{EV1 --- Constant-demand convergence}
\label{sec:meth:ev1}

A constant daily demand of $d$ units is applied to all retailers over
a 100-day horizon.  Each node runs a base-stock policy with
$S = d(L+1)$ where $L$ is the lead time from its parent.  The test
confirms that on-hand inventory at each node converges to and stays at
$S - d \cdot L$ (the equilibrium on-hand level) within $L+1$ days.

\subsection{EV2 --- Demand shock response}
\label{sec:meth:ev2}

A synthetic demand series with constant background demand $d$ and a
single large spike at day 20 is read from \texttt{inputs/demand\_shock.csv}.
The experiment verifies that the upstream nodes observe an amplified order
spike with a delay matching the cumulative pipeline lead time.

\subsection{EV3 --- Policy comparison under transport capacity}
\label{sec:meth:ev3}

Two policies (base-stock and $(s,S)$) are compared on the single-SKU
topology with and without a finite vehicle capacity.  With no capacity
limit both policies track demand similarly.  With capacity limits, the
$(s,S)$ fixed-quantity orders align better with vehicle loads.

\subsection{EV4 --- Capacity scaling}
\label{sec:meth:ev4}

Vehicle capacity is scaled by multipliers $\{0.25, 0.5, 1.0, 2.0, 4.0\}$
relative to a baseline calibrated to carry one week of average demand.
M5 Walmart demand data is used for realism.  Results appear in
Table~\ref{tab:ev4} (Section~\ref{sec:results:ev4}).

%% -----------------------------------------------------------------------
\section{Dataset: M5 Walmart Forecasting Data}
\label{sec:meth:dataset}

The demand data for all experiments is derived from the
\textbf{M5 Walmart Forecasting dataset} \citep{MakridakisEtAl2020}.
The M5 dataset contains daily unit sales for thousands of products across
ten Walmart stores in the United States, spanning approximately five years.

The M5 dataset was selected because it exhibits realistic retail demand
characteristics that stress inventory policies beyond what synthetic
Poisson processes can capture:
\begin{itemize}
  \item \textbf{Intermittency and spikes.} Sales often cluster around
        promotional events and holidays, creating irregular demand patterns.
  \item \textbf{Weak weekly seasonality.} Weekend peaks and weekday troughs
        introduce low-frequency variation.
  \item \textbf{Long tails.} Occasional large orders produce high
        coefficient-of-variation values that challenge base-stock stability.
\end{itemize}

Per-retailer daily demand CSVs (\texttt{inputs/demand\_data/R1.csv},
\texttt{R2.csv}, \texttt{R3.csv}) were derived by selecting one M5
series per retailer per SKU, filtering to a 365-day window, and reindexing
to a zero-based day index.  For the preliminary validation experiments,
a separate synthetic series (\texttt{demand\_shock.csv}) with a single
spike at day 20 was constructed.

%% -----------------------------------------------------------------------
\section{Main Experiment Configuration}
\label{sec:meth:config}

All main experiments (E0--E4) use the configuration file
\texttt{config/1n3\_5sku.json} with the following parameters:

\begin{table}[htbp]
  \centering
  \caption{Main experiment simulation parameters.}
  \label{tab:sim-params}
  \begin{tabular}{@{}ll@{}}
    \toprule
    Parameter & Value \\
    \midrule
    Network topology & 1 supplier, 1 warehouse (W1), 3 retailers (R1--R3) \\
    SKUs & 5 (SKU\_A through SKU\_E) \\
    Simulation horizon & 365 days \\
    Warm-up period & 30 days (excluded from KPIs) \\
    Demand type & CSV-driven (M5 Walmart, retailer-specific) \\
    Lead times & Deterministic (supplier$\to$W1: 2d; W1$\to$Ri: 1d) \\
    Fill-rate target & $\geq$ 92\% (soft constraint, experiments E1b--E3) \\
    \bottomrule
  \end{tabular}
\end{table}

SKU physical properties span a range of volumes (0.002--0.008 m$^3$/unit)
and weights (0.5--1.2 kg/unit), so the GreedyLoadPlanner's dual
volume-and-weight capacity constraint is active for some SKU mixes.

%% -----------------------------------------------------------------------
\section{Warm-Up Exclusion}
\label{sec:meth:warmup}

At simulation start, all nodes are initialised with a configurable amount
of stock.  This initial stock creates a transient period during which
KPIs reflect the initialisation state rather than the steady-state
behaviour of the policy.  To avoid this bias, the first 30 days of each
simulation run are \emph{excluded} from all reported KPIs (costs, fill
rates, bullwhip ratios).

Warm-up exclusion is enforced by passing the \texttt{--warmup 30} flag
to all experiment scripts.  The warm-up period was chosen to be
$\geq 2 \times L_{\max}$, where $L_{\max} = 2$ days is the maximum
lead time in the network, ensuring that initial inventory effects have
dissipated before measurement begins.

%% -----------------------------------------------------------------------
\section{Multi-Seed Statistical Validation}
\label{sec:meth:multiseed}

To quantify the sensitivity of results to the random seed (which affects
stochastic demand realisations and TPE search paths), the framework
provides \texttt{scripts/multi\_seed\_validation.py}.  This script:

\begin{enumerate}
  \item Runs the simulation with the specified policy parameters across
        $N$ independently seeded runs (default $N = 10$).
  \item Reports the 95\% confidence interval on total cost and fill rate
        using the normal approximation.
  \item Compares two configurations (e.g.\ E0 vs E3) using a
        Welch two-sample $t$-test to assess whether the observed
        cost difference is statistically significant.
\end{enumerate}

Because the main experiments use a fixed deterministic demand CSV,
cross-seed variance in cost is primarily driven by the optimiser's
stochastic search path rather than demand variability.  Multi-seed
validation is most informative when used with Poisson demand generators.

%% -----------------------------------------------------------------------
\section{Bullwhip Metric}
\label{sec:meth:bullwhip}

The bullwhip ratio is the ratio of the order coefficient-of-variation
squared to the demand coefficient-of-variation squared:

\begin{equation}
  \text{BW}_{n,s} = \frac{\text{CV}^2(\text{orders}_{n,s})}
                         {\text{CV}^2(\text{demand}_{n,s})}
  \label{eq:bullwhip}
\end{equation}

A ratio greater than one at node $n$ indicates that orders placed by $n$
are \emph{more} variable than the demand $n$ faces, i.e.\ amplification
is occurring.  The metric is computed on the post-warm-up window only,
using the daily order logs stored in \texttt{orders\_log.csv}.

For retailers, the denominator is the external customer demand; for the
warehouse and supplier, the denominator is the aggregate demand from
their children (which is itself the retailers' orders).

%% -----------------------------------------------------------------------
\section{Main Optimisation Experiments E0--E4}
\label{sec:meth:main}

The six main experiments are summarised in Table~\ref{tab:exp-overview}.

\begin{table}[htbp]
  \centering
  \caption{Main optimisation experiments: design and hypothesis.}
  \label{tab:exp-overview}
  \renewcommand{\arraystretch}{1.2}
  \begin{tabular}{@{}llp{7cm}@{}}
    \toprule
    ID & Title & Hypothesis \\
    \midrule
    E0 & Analytical baseline &
         Newsvendor stock levels + zero dispatch thresholds will produce
         a high transport cost and sub-92\% fill rate. \\[4pt]
    E1a & Fixed 25\% dispatch threshold &
          Imposing a uniform consolidation constraint without adjusting
          inventory will reduce transport cost but collapse fill rate. \\[4pt]
    E1b & Transport-only optimisation &
          Optimising dispatch thresholds alone (inventory fixed at E0)
          will achieve modest cost reduction but cannot reach 92\% fill. \\[4pt]
    E2 & Inventory-only optimisation &
         Optimising stock levels alone (thresholds fixed at 0) will
         achieve 92\% fill but increase total cost. \\[4pt]
    E3 & \textbf{Joint optimisation (main result)} &
         Simultaneous optimisation of inventory and transport will
         achieve the lowest cost at $\geq$92\% fill. \\[4pt]
    E4 & Pareto frontier &
         The cost--service trade-off exhibits diminishing returns: each
         percentage point of fill rate above 95\% costs substantially more. \\
    \bottomrule
  \end{tabular}
\end{table}

All experiments except E0 use the Optuna-TPE optimiser described in
Chapter~\ref{chap:optimization}.  Experiment E0 uses analytically derived
parameters and requires no optimisation.  All experiments share the same
simulator configuration, demand data, and warm-up window to enable fair
comparison.

%% -----------------------------------------------------------------------
\section{Reproducibility}
\label{sec:meth:reproducibility}

Each experiment run generates a self-contained output directory at
\texttt{outputs/experiments\_\textlangle ts\textrangle /}, containing:

\begin{itemize}
  \item \texttt{params.json} — the exact policy parameters used.
  \item \texttt{summary.txt} — total cost, cost breakdown, fill rates,
        and bullwhip ratios in human-readable form.
  \item CSV logs: \texttt{inventory\_log.csv}, \texttt{orders\_log.csv},
        \texttt{costs\_by\_node\_sku.csv}, \texttt{kpis\_by\_sku.csv},
        \texttt{bullwhip.csv}, \texttt{shipments\_log.csv}.
  \item PNG plots: cost breakdown, fill rate, inventory time-series,
        order time-series, bullwhip, and Pareto frontier.
\end{itemize}

The CLI commands to reproduce all results are given in
Appendix~\ref{appn:reproduction}.
